\subsection{Environment}
\begin{itemize}[leftmargin=*]
    \item \textbf{Households:} overlapping generations with two life stages: working and retired. Incomplete markets. Idiosyncratic risk: income and employment. Survival is stochastic and age-dependent.
    \item \textbf{Firms:} representative firm with Cobb--Douglas technology using private and public capital.
    \item \textbf{Government:} sets tax rates $(\tau^c,\tau^l,\tau^p,\tau^k)$, pension rules, health coverage rate $\kappa$, government consumption $G_t$, and public investment $I_t^g$. Issues debt $B_t$.
    \item \textbf{Interest rate:} exogenous path $\{r_t\}$. Given $r_t$, the firm's first-order conditions determine the wage $w_t$ and the private capital-labor ratio.
\end{itemize}

\subsection{Demographics and Life Cycle}
Ages $j \in \{1,\dots,J\}$:
\begin{align*}
\text{Working: } & j \in \mathcal{W}=\{1,\dots,J_R\}, \\
\text{Retired: } & j \in \mathcal{R}=\{J_R+1,\dots,J\}.
\end{align*}
Survival probability $\pi(j)$ depends on age. Population grows at exogenous rate $g_N$; cohort sizes are weighted by $\Lambda^{-j}$ where $\Lambda = 1+g_N$.

\subsection{Preferences}
Period utility:
\begin{equation}
u(c,\ell) = \frac{c^{1-\sigma}}{1-\sigma} - \nu\,\frac{\ell^{1+\varphi}}{1+\varphi},
\end{equation}
where $c$ is consumption and $\ell \in [0,1]$ is hours worked.

\subsection{Health Expenditure}
Each household faces a fixed per-period medical expenditure $m$ (expressed as a fraction of mean income). The government covers fraction $\kappa$ of this cost; the household pays the remainder $(1-\kappa)m$ out of pocket. Aggregate government health spending is $\text{HSub}_t = \kappa\,m\,N_t$, where $N_t$ is total population.

$m$ is calibrated to match total health expenditure as a share of GDP, taken from national accounts (Eurostat/OECD Health Statistics). $\kappa$ is set to the public share of total health expenditure in the calibrated country.

\subsection{Idiosyncratic Income and Employment Risk}
Income state $z \in \{0, z_1, \dots, z_{n_y-1}\}$ follows a Markov chain $P_z(z'|z)$, where $z=0$ represents unemployment. The process is discretized from a log-normal AR(1) using Tauchen's method. Wage income during working years:
\begin{equation}
y^L = w_t\,\kappa_j\,z,
\end{equation}
where $\kappa_j$ is a deterministic age productivity profile. When $z=0$, the household receives UI benefits in place of labor income.

\subsection{Household Budget Constraints}\label{sec:household-bc}

\subsubsection*{Working ($j \in \mathcal{W}$)}
\begin{equation}\label{eq:bc-work}
(1+\tau^c)\,c + a' = \bigl(1+(1-\tau^k)r_t\bigr)\,a + (1-\tau^l-\tau^p)\,y^L\,\ell + T^{UI} - (1-\kappa)\,m,
\end{equation}
with $a' \geq \underline{a}$. The UI transfer is $T^{UI} = \rho^{ui}\,w_t\,\kappa_j\,z_{\text{last}}$ when $z=0$, where $z_{\text{last}}$ is the last-employed income state. A minimum consumption floor is enforced via a lump-sum transfer when needed.

\subsubsection*{Retired ($j \in \mathcal{R}$)}
\begin{equation}\label{eq:bc-ret}
(1+\tau^c)\,c + a' = \bigl(1+(1-\tau^k)r_t\bigr)\,a + (1-\tau^l)\,\text{PENS} - (1-\kappa)\,m,
\end{equation}
with $a' \geq \underline{a}$. The pension benefit is:
\begin{equation}\label{eq:pension}
\text{PENS} = \max\!\left(\rho\,w_t\,\kappa_{J_R}\bigl[\lambda\, z_{\text{last}} + (1-\lambda)\,\bar{z}\bigr],\; b_{\min}\right),
\end{equation}
where $\rho$ is the replacement rate, $\lambda$ weights last-period income against mean employed income $\bar{z}$, and $b_{\min}$ is a pension floor.

\subsubsection*{Terminal condition}
Assets of the deceased are taxed at rate $\tau^{\text{beq}}$ and redistributed lump-sum to surviving households.

\subsection{Production}
A representative firm produces with Cobb--Douglas technology:
\begin{equation}\label{eq:production}
Y_t = A\,\bigl(K_t^g\bigr)^{\eta_g}\,K_t^{\alpha}\,L_t^{1-\alpha},
\end{equation}
where $K_t$ is private capital, $K_t^g$ is public capital, $L_t$ is aggregate effective labor, $\alpha$ is the private capital share, $\eta_g$ is the output elasticity of public capital, and $A$ is TFP. Public capital evolves as:
\begin{equation}\label{eq:kg}
K_{t+1}^g = (1-\delta_g)\,K_t^g + I_t^g.
\end{equation}
Given exogenous $r_t$, the firm's first-order conditions determine:
\begin{align}
\frac{K_t}{L_t} &= \left(\frac{\alpha A\,(K_t^g)^{\eta_g}}{r_t+\delta}\right)^{1/(1-\alpha)}, \label{eq:foc-k}\\
w_t &= (1-\alpha)\,A\,(K_t^g)^{\eta_g}\left(\frac{K_t}{L_t}\right)^{\alpha}. \label{eq:foc-w}
\end{align}

\subsection{Individual Problem}\label{sec:state}
Individual state: $(j,\, z,\, z_{\text{last}},\, a)$, where $z_{\text{last}}$ is the last-employed income state, frozen during unemployment and updated to $z$ when $z>0$.

Aggregate state: $\mathbf{S}_t = (r_t, w_t, K_t^g, \mu_t)$, where $\mu_t$ is the cross-sectional distribution.

\paragraph{Working ($j \in \mathcal{W}$).}
\begin{align*}
V_j(z, z_{\text{last}}, a;\,\mathbf{S}) = \max_{c,\,\ell,\,a' \geq \underline{a}} \;& u(c,\ell) + \beta\,\pi(j)\,\mathbb{E}\bigl[V_{j+1}(z', z_{\text{last}}', a';\,\mathbf{S}')\bigr] \\
\text{s.t. } & \eqref{eq:bc-work}, \quad z' \sim P_z(\cdot|z), \\
& z_{\text{last}}' = z \cdot \mathbf{1}[z>0] + z_{\text{last}} \cdot \mathbf{1}[z=0].
\end{align*}

\paragraph{Retired ($j \in \mathcal{R}$).}
\begin{align*}
V_j^R(z_{\text{last}}, a;\,\mathbf{S}) = \max_{c,\,a' \geq \underline{a}} \;& u(c,0) + \beta\,\pi(j)\,\mathbb{E}\bigl[V_{j+1}^R(z_{\text{last}}, a';\,\mathbf{S}')\bigr] \\
\text{s.t. } & \eqref{eq:bc-ret}.
\end{align*}

\subsection{Taxes, Pensions, and Transfers}
\begin{itemize}[leftmargin=*]
    \item \textbf{Taxes:} four proportional rates $(\tau^c, \tau^l, \tau^p, \tau^k)$ on consumption, labor income, payroll, and capital income. An optional progressive schedule replaces the flat labor rate: $\tau^l(y) = 1 - \kappa^{\text{tax}}\,y^{-\eta^{\text{tax}}}$.
    \item \textbf{Pensions:} replacement-rate formula \eqref{eq:pension}, pay-as-you-go. Payroll tax revenues feed an optional pension trust fund: $S_{t+1}^{\text{pens}} = (1+r_t)\,S_t^{\text{pens}} + \text{Rev}_t^p - \text{PENS}_t^{\text{out}}$.
    \item \textbf{UI:} replacement rate $\rho^{ui}$ applied to $z_{\text{last}}$.
    \item \textbf{Transfer floor:} lump-sum transfer enforcing a minimum consumption level.
    \item \textbf{Bequests:} assets of the deceased are taxed at rate $\tau^{\text{beq}}$ and redistributed uniformly.
\end{itemize}

\subsection{Government Budget}\label{sec:gov-budget}
Government revenues:
\begin{align*}
\text{Rev}_t^c &= \tau_t^c\,C_t, \quad
\text{Rev}_t^l = \tau_t^l\,\mathcal{B}_t^{\text{lab}}, \quad
\text{Rev}_t^p = \tau_t^p\,\mathcal{B}_t^{\text{lab}}, \\
\text{Rev}_t^k &= \tau_t^k\,r_t\,A_t, \quad
\text{BeqTax}_t = \tau^{\text{beq}}\int (1-\pi(j))\,a\,d\mu_t,
\end{align*}
where $\mathcal{B}_t^{\text{lab}} \equiv \int_{j\in\mathcal{W}} w_t\,\kappa_j\,z\,\ell\,d\mu_t$ and $A_t \equiv \int a\,d\mu_t$.

Spending aggregates:
\begin{align*}
\text{UI}_t &= \int_{j\in\mathcal{W},\,z=0} T^{UI}\,d\mu_t, \quad
\text{PENS}_t^{\text{out}} = \int_{j\in\mathcal{R}} \text{PENS}\,d\mu_t, \quad
\text{HSub}_t = \kappa\,m\,N_t.
\end{align*}

Primary deficit and debt dynamics:
\begin{align}
\text{PD}_t &= \text{UI}_t + \text{PENS}_t^{\text{out}} + \text{HSub}_t + G_t + I_t^g - \bigl(\text{Rev}_t^c + \text{Rev}_t^l + \text{Rev}_t^p + \text{Rev}_t^k + \text{BeqTax}_t\bigr), \label{eq:primary-deficit}\\
B_{t+1} &= (1+r_t)\,B_t + \text{PD}_t. \label{eq:debt}
\end{align}

\subsection{Aggregation and Market Clearing}
Aggregate labor and output:
\begin{align*}
L_t &= \int_{j\in\mathcal{W}} \kappa_j\,z\,\ell\,d\mu_t, \\
Y_t &= A\,(K_t^g)^{\eta_g}\,K_t^\alpha\,L_t^{1-\alpha}.
\end{align*}
Given exogenous $r_t$, the firm's first-order conditions \eqref{eq:foc-k}--\eqref{eq:foc-w} pin $K_t/L_t$ and $w_t$; private capital follows as $K_t = (K_t/L_t)\cdot L_t$. Household wealth $A_t = \int a\,d\mu_t$ need not equal $K_t + B_t$: the gap is net foreign assets, $\text{NFA}_t = A_t - K_t - B_t$. Debt accumulates through \eqref{eq:debt}.

The economy is simulated via Monte Carlo. At each period, $n_{\text{sim}}$ households per cohort-education cell draw idiosyncratic shocks and follow optimal decision rules. Aggregates are constructed by summing across cohorts weighted by $\Lambda^{-j}$.

\subsection{Model Features Not Yet Implemented}\label{sec:not-implemented}

\begin{enumerate}
    \item \textbf{Age- and income-bracket heterogeneity in medical expenditure.} The fixed cost $m$ is uniform across agents. Medical expenditure rises with age and falls with income in the data. Extending the model to $m(j, y)$ would allow the health spending channel to interact with the income distribution and the age structure of the population.
    \item \textbf{Schooling phase and child costs.} The theory includes child consumption costs $c_y(j)$ and private education expenditure $e(j) = (1-\kappa^{\text{school}}_t)\bar{e}(j)$ during working years. Neither is in the code. See Appendix~\ref{s:calibration} for the associated data requirements.
    \item \textbf{Robustness.} The Cobb--Douglas production function \eqref{eq:production} can be replaced by the CES aggregator $\tilde{K} = \bigl[(1-\theta)K^\rho + \theta(K^g)^\rho\bigr]^{1/\rho}$, which nests Cobb--Douglas at $\rho \to 0$. We use Cobb--Douglas as the baseline and consider the CES specification as a robustness check.
    \item \textbf{Pension rule.} The formula in equation~\eqref{eq:pension} is a parsimonious approximation. The pension rule for the calibrated country should be checked against national legislation and the OECD Pensions at a Glance database; the replacement rate $\rho$, the earnings base, and the floor $b_{\min}$ will be adjusted to match the country-specific rule at the calibration stage.
\end{enumerate}




